\section{Baltic}

The \textbf{Baltic} tradition of the Lithuanian and Latvian peoples is reconstructed and fragmentary. No native texts survive from the pagan period, and knowledge derives from Roman and German chronicles, from etymology, from comparative reconstruction, and above all from later folk songs, the Latvian \textbf{dainas}. The pantheon is conservative and preserves an old Indo-European pattern.

\subsection{The Creation Model}

There is no single preserved account of the making of the cosmos. What survives is an ordering of the sky rather than a told beginning. At the head stands a sky father, the supreme principle of the heavens, who fixes the cosmic order. From him come the divine twins, the sons of the sky.

The visible world is governed by a celestial family. A solar principle, the sun maiden, is held to be responsible for all life on earth and to carry the souls of the dead across the sea, dying and being reborn with each day. A lunar principle, the moon, is set against her as her partner. In the songs the moon is unfaithful and is punished by the thunder principle, which splits him, and so the day and the night are divided. An axis tree stands on the path of the sun, by the water at the edge of the world, joining the realms of the sky. Around it turns a recurring celestial wedding, in which the daughter of the sun is courted by the sons of the sky, by the morning star, and by the thunderer.

\subsection{Key Motifs}

The sky father is \textbf{Dievas}. The sun goddess is \textbf{Saule}, with her daughter the \textbf{Saules meita}. The moon is \textbf{Menulis}, in Latvian \textbf{Meness}, and the thunder god is \textbf{Perkunas}, in Latvian \textbf{Perkons}, who punishes the moon for his infidelity with the morning star \textbf{Ausrine}. The world tree is the \textbf{Austras koks}, the tree of dawn, standing on the path of the sun by the water. The midsummer solstice festival is the Lithuanian \textbf{Rasos} and the Latvian \textbf{Ligo}.

\begin{longtable}{@{}
ll@{}}
\toprule
Lithuanian & Role \\
\midrule
\textbf{Dievas} & sky father, supreme order \\
\textbf{Saule} & the sun, source of life \\
\textbf{Menulis} & the moon, partner of the sun \\
\textbf{Perkunas} & the thunder, keeper of order \\
\bottomrule
\end{longtable}

The Baltic material fits the shared pattern of the one going out into the many, since a single sky father gives rise to the celestial family of sun, moon, and thunder whose ordered turning sustains the world around the axis tree.
